\documentclass{ximera}
\title{Exponential Growth Summary}

\newcommand{\pskip}{\vskip 0.1 in}

\begin{document}
\begin{abstract}
Exponential functions.
\end{abstract}
\maketitle


\section{Growth Factors}
A population that grows \emph{linearly} increases at a constant absolute rate, perhaps at the rate of $20,000$ people/yr. 

A population that grows \emph{exponentially} increases at a constant \emph{relative} rate, perhaps at the rate of $20\%$/yr. To understand what this means, suppose the population was $P$ at the start of the year 2020. Then one year later the population would be $20\%$ greater, or 
\[
   P + 20\% P  = P + 0.2 P = 1.2P .
\]
We call 1.2 the \emph{annual growth factor}. To get the population one year from, we just multiply the current population by this factor of 1.2.


\begin{example} \label{Ex8mvemaz}

Between 12:15pm and 10pm the population of a colony of bacteria increases exponentially. The population is $300,000$ at 2pm and $360,000$ at 3pm.

\begin{enumerate}
\item Describe how the population grows.

\item Find the one-hour growth factor and explain its meaning.

\item Find the population at 10pm.

\item Find a function 
\[
     P=f(t) \, \, \, , 1/4\leq t \leq 10, 
\]
that expresses the population in terms of the number of hours past noon.

\item Find the population one hour after there are 400,000 bacteria.

\item Find the population one hour before there are 500,000 bacteria.

\item Find the population 90 minutes after there are 500,000 bacteria.
\end{enumerate} 

\begin{explanation}
\begin{enumerate}
\item Between 2pm and 3pm the population increases by $60,000$ bacteria. So every hour the population increases by
\[
     \frac{60,000\text{ bac}}{300,000\text{ bac}} = 20\%.
\]

\item The one-hour growth factor is
\[
 \frac{360,000\text{ bac}}{300,000\text{ bac}} = 1.2 .
\]
So to get the popluation one hour from now, we multiply the current population by $1.2$.

\item There are $10-2 = 8$ hours from 2pm to 10pm. So to get the population at 10pm, we multiply the population at 2pm by the one-hour growth factor eight times. This means the population at 10pm is
\[
   300,000(1.2)^{10-2} = 300,000(1.2)^8.
\]

\item The same reasoning as in part (c) shows that
\begin{align*}
 P &= f(t) \\
    &= f(2) (1.2)^{t-2} \\
    &= 300,000(1.2)^{t-2} \, , \, 1/4\leq t \leq 10 .
\end{align*}

Or that
\begin{align*}
 P &= f(t) \\
    &= f(3) (1.2)^{t-3} \\
    &= 360,000(1.2)^{t-3} \, , \, 1/4\leq t \leq 10.
\end{align*} 

Note: It is not necessary, nor is it proper, to find the population at noon. We cannot determine the population at this time from the given information. 

\item It is best not to use our function for this. We use the one-hour growth factor instead. Then one-hour after there are 400,000 bacteria, the population is
\[
  (400,000 \text{ bac})(1.2) = 480,000 \text{ bac}.
\]

\item To find the population one-hour before there are 500,000, we just divide 500,000 by the one-hour growth factor to get approximately
\[
 500,000(1.2)^{-1} = \frac{500,000}{1.2} \sim 416,667
\]
bacteria.

\item We first find the $30$-minute growth factor. This is \emph{not} 1.1 because even though the population increases by $20\%$ every hour it does \emph{not} increase by $10\%$ every half-hour. For if the population did increase by $10\%$ every half-hour, the one-hour growth factor would be
\[
    1.1^2 = 1.21
\]
and not 1.2 (ie. the population would increase by $21\%$ every hour if it increased by $10\%$ every half-hour).

We'll let $g$ be the half-hour growth factor. Then the one-hour growth factor would be $g^2$, so
\[
    g^2 = 1.2 .
\]
And since $g>0$, 
\[
  g = \sqrt{1.2} = 1.2^{1/2} \sim 1.095445 .
\] 

Then the $90$-minute growth factor is
\[
 g^3 = \left( 1.2^{1/2} \right)^3 = 1.2^{3/2}.
\]

So $90$ minutes after there are 500,000 bacteria, the population is
\[
  (500,000 \text{ bac}) \left(  1.2 \right)^{3/2} \sim 657,267 \text{ bac}.
\]
\end{enumerate}
\end{explanation}

\end{example}


\begin{example} \label{Ex8femvemaz}

Between 12:15pm and 8pm the population of a colony of bacteria increases exponentially. The population is $300,000$ at 2pm and $350,000$ at 5pm.

Find a function 
\[
  P=f(t) \, \, \, , 1/4\leq t \leq 10,
\]
that expresses the population in terms of the number of hours past noon.

\begin{explanation}
The key is to find the one-hour growth factor. 

So we let $g$ be the one-hour growth factor. Then $g^3$ is the three-hour growth factor.

But the three-hour growth factor is
\[
  \frac{350,000\text{ bac}}{300,000\text{ bac}} = \frac{7}{6}.
\]

So
\[
  g^3 = 7/6
\] 
and
\[
  g = \left( \frac{7}{6} \right)^{1/3}.
\]

So
\begin{align*}
     P  &= f(t) \\
         &= f(2) g^{t-2} \\
         &= 300,000 \left(   \left( \frac{7}{6} \right)^{1/3}\right)^{t-2} \\
         &= 300,000 \left( \frac{7}{6} \right)^{\frac{1}{3}\left( t-2 \right)}.
\end{align*}

Or,
\begin{align*}
     P  &= f(t) \\
         &= f(5) g^{t-\answer{5}} \\
         &= \answer{350000} \left(   \left( \frac{7}{6} \right)^{1/3}\right)^{t-\answer{5}} \\
         &= \answer{350,000} \left( \frac{7}{6} \right)^{\frac{1}{3}\left( t-\answer{5} \right)}.
\end{align*}

\end{explanation}

\end{example}

\end{document}
